\documentclass[11pt,a4paper]{article}

% --- PACOTES ESSENCIAIS ---
\usepackage[utf8]{inputenc}
\usepackage[T1]{fontenc}
\usepackage[portuguese]{babel}
\usepackage[a4paper,margin=2.5cm]{geometry}
\usepackage{graphicx}
\usepackage[table]{xcolor}
\usepackage{titlesec}
\usepackage{booktabs}
\usepackage{tabularx}
\usepackage{enumitem}
\usepackage{hyperref}
\usepackage{setspace}
\usepackage{longtable}
\usepackage{float}

% --- DEFINIÇÃO DE CORES ---
\definecolor{petrolblue}{RGB}{15,50,90}
\definecolor{petrolgold}{RGB}{200,160,50}
\definecolor{darkgray}{RGB}{50,50,50}
\definecolor{lightgray}{RGB}{245,247,250}

% --- CONFIGURAÇÃO DE LINKS ---
\hypersetup{colorlinks=true,linkcolor=petrolblue,urlcolor=petrolblue}

% --- FORMATAÇÃO DE TÍTULOS ---
\titleformat{\section}{\normalfont\Large\bfseries\color{petrolblue}}{\thesection}{1em}{}{\vspace{0.2em}}\titleformat{\subsection}{\normalfont\large\bfseries\color{darkgray}}{\thesubsection}{1em}{}
\setlength{\parskip}{0.5em}
\setlength{\parindent}{0pt}

\begin{document}

\begin{titlepage}
    \centering
    \vspace*{1.8cm}

    {\Huge\textbf{\color{petrolblue} PETROLÚMEN}}\\[0.5cm]
    {\Large\textbf{Sistema Integrado de Engenharia de Reservatórios e Previsão Preditiva}}\\[1.3cm]

    {\large\textit{“Transformar a análise de dados de petróleo e gás: substituir o caos do Excel por um ecossistema desktop inteligente, seguro e acessível.”}}\\[2.4cm]

    {\Large\textbf{Memorandum Estratégico e de Viabilidade Técnico-Financeira}}\\[1.8cm]

    \begin{tabular}{rl}
        \textbf{Versão:} & 1.0 — Candidatura GEOENERGY 2026 \\
        \textbf{Data:} & Outubro de 2026 \\
        \textbf{Âmbito:} & Angola e mercado regional de energia \\
        \textbf{Categoria:} & Digitalização e Tecnologias Aplicadas às Geociências e Energia \\
        \textbf{Classificação:} & Documento de submissão \\
    \end{tabular}

    \vfill

    {\large\textbf{Equipa:}}\\
    Rocélio Da Silva (Proponente Principal)\\
    Derik Francisco\\
    Davidaldo Belchior\\
    Mafra Tito\\
    \textit{Estudantes de Engenharia de Petróleos — ISPTEC}

    \vspace{1cm}
\end{titlepage}

\tableofcontents
\newpage

\section{Sumário Executivo}
A indústria de petróleo e gás em Angola e na África Subsaariana enfrenta um gargalo tecnológico silencioso: a maior parte do trabalho de engenharia de reservatórios é ainda apoiada por folhas de cálculo manual, ferramentas fragmentadas e software empresarial de elevado custo. O resultado é perda de tempo, risco de erros humanos e baixa capacidade de aplicar modelos preditivos em operações do dia a dia.

O \textbf{PetroLúmen} surge como resposta a este problema. Trata-se de uma plataforma desktop de engenharia de reservatórios, concebida para funcionar localmente, com arquitetura integrada em \textit{frontend} React/Next.js e \textit{backend} Python/FastAPI, apoiada por Tauri para distribuição nativa. O produto reúne PVT, análise de declínio, previsão de produção, simulação técnica e inteligência artificial num único ecossistema, com foco em usabilidade, segurança e acessibilidade.

A análise do repositório confirma que a base tecnológica do produto já está materializada. O backend inclui módulos de PVT, análise de reservatórios, simulação, histórico, declínio e previsão; o frontend é uma aplicação web empacotada para desktop; e os testes automatizados do backend são um indicador claro de maturidade técnica. Assim, o valor do projeto em contexto GEOENERGY não está apenas na ideia, mas na combinação de solução realista, engenharia validada e potencial de implementação no mercado angolano e regional.

\section{Análise do Repositório e Evidência Técnica}
A análise técnica do código revelou uma estrutura coerente e orientada para produto industrial:

\begin{itemize}
    \item \textbf{Frontend:} diretório \texttt{petrolumen/}, com aplicação Next.js, componentes UI e integração Tauri para desktop. O \texttt{package.json} confirma o uso de React, Next, Tailwind, shadcn/ui, Vitest e Tauri.
    \item \textbf{Backend:} diretório \texttt{backend/}, com FastAPI, routers da API e uma estrutura modular baseada em \texttt{gaia	extunderscore genesis}. Os módulos incluem PVT, simulação, análise de produção, previsão, otimização, RTA (well testing), geomecânica e\ analysis.
    \item \textbf{Testes automatizados:} a pasta \texttt{backend/tests/} inclui testes de auth, PVT, análise, simulação, background jobs e concorrentemente, mostrando um esforço real de validação funcional.
    \item \textbf{Capacidade de produto:} o repositório contempla fluxos de upload de dados, estabilização de sessão, processamento assíncrono, análise de declínio, previsão com IA e documentação técnica organizada em \texttt{README.md}, \texttt{docs/} e \texttt{backend/docs	extunderscore sphinx/}.
\end{itemize}

Em termos de maturidade, o projeto apresenta uma base tecnológica suficiente para ser tratado como prova de conceito com potencial de desenvolvimento comercial em curto prazo, especialmente para um concurso académico com foco em produto e empreendedorismo.

\section{O Problema e a Oportunidade Estratégica}
\subsection{O "Pesadelo do Excel" e a falta de acesso a software especializado}
\begin{itemize}
    \item Engenheiros de reservatórios e estudantes perdem horas trabalhando com folhas Excel e cálculos dispersos.
    \item Ferramentas enterprise são caras, pesadas e pouco acessíveis para pequenas empresas, consultores independentes e universidades.
    \item A ausência de um ambiente integrado impede a adoção de modelos de previsão, análise de declínio e simulação local em contextos reais.
\end{itemize}

\subsection{A oportunidade}
O mercado angolano e regional tem um espaço específico para uma solução acessível, local e orientada ao desempenho técnico. O PetroLúmen combina três elementos estratégicos: 
\begin{enumerate}
    \item software local e seguro, sem dependência de nuvem pública;
    \item cálculo técnico robusto para PVT, DCA, RTA e simulação;
    \item modelo de negócio escalável, começando por estudantes e consultores e evoluindo para empresas.
\end{enumerate}

\section{Visão de Produto e Modelo de Negócio}
\subsection{Product Canvas do PetroLúmen}
\begin{table}[H]
\centering
\renewcommand{\arraystretch}{1.25}
\begin{tabularx}{\textwidth}{|l|X|}
\hline
\rowcolor{petrolblue}\textcolor{white}{\textbf{Componente}} & \textcolor{white}{\textbf{Descrição}} \\
\hline
\textbf{Para} & Engenheiros de reservatórios, estudantes, consultores independentes e pequenas operadoras. \\
\hline
\textbf{Cujo problema} & A análise técnica é feita em ferramentas fragmentadas, with Excel ou software caro e inacessível. \\
\hline
\textbf{O produto} & PetroLúmen é uma plataforma desktop de engenharia de reservatórios, com módulos de PVT, análise de declínio, simulação, previsões e IA. \\
\hline
\textbf{Que oferece} & Segurança local, acessibilidade, rapidez e consistência nos cálculos. \\
\hline
\textbf{Diferenciação} & Arquitetura local, pricing acessível, integração de IA e foco em uso diário de reservatórios. \\
\hline
\textbf{Modelo de negócio} & Freemium para estudantes, Pro para consultores e Enterprise para empresas. \\
\hline
\end{tabularx}
\end{table}

\subsection{Arquitetura funcional proposta}
\begin{enumerate}
    \item \textbf{Módulo PVT:} cálculo de fatores de formação, viscosidade, compressibilidade e correlações empíricas.
    \item \textbf{Módulo de Declínio:} modelos Arps, Duong, PLE e SEPD, com ajuste e previsão.
    \item \textbf{Módulo de Simulação:} motores de reservatório, balance de materiais e análise de fluxo.
    \item \textbf{Módulo de IA:} modelos para previsão de produção e séries temporais.
    \item \textbf{Módulo de Gestão e Segurança:} dados locais, sessões isoladas e interface desktop segura.
\end{enumerate}

\section{Mercado, Clientes-Alvo e Posicionamento}
O PetroLúmen não pretende competir diretamente com plataformas 3D e geológicos de grande escala, mas sim com o problema operativo do dia a dia. O foco é claro:

\begin{itemize}
    \item \textbf{Mercado primário (B2C):} estudantes de Engenharia de Petróleos, investigadores e jovens profissionais.
    \item \textbf{Mercado secundário (B2B):} consultores independentes, pequenas consultoras e operadores de campos marginais.
    \item \textbf{Mercado terciário (Enterprise):} empresas nacionais e regionais com necessidade de segurança local e personalização.
\end{itemize}

A estratégia é iniciar no mercado académico e consultivo, criar um ciclo de adoção e, depois, escalar para a operação empresarial. Esta abordagem reduz o risco, permite validação técnica e protege a eficiência do modelo de receitas.

\section{Viabilidade Técnica e Recursos Necessários}
\begin{itemize}
    \item \textbf{Tecnologias base:} Python (FastAPI), React.js, Next.js, Tauri e SQLite/PostgreSQL.
    \item \textbf{Estado do repositório:} arquitetura e módulos já implementados em várias áreas da engenharia de reservatórios.
    \item \textbf{Validação:} a pasta \texttt{backend/tests/} evidencia um esforço real de validação funcional, com foco em endpoints, cálculo, concorrência e simulações.
    \item \textbf{Recursos da equipa:} dois estudantes de Engenharia de Petróleos do ISPTEC, com perfil técnico complementar e capacidade de desenvolvimento aplicando engenharia e empreendedorismo.
\end{itemize}

\section{Modelo de Receitas}
O modelo de negócio do PetroLúmen assenta em três níveis:
\begin{enumerate}
    \item \textbf{Freemium (estudantes):} acesso gratuito com limitação de exportação e uso de alguns módulos.
    \item \textbf{Pro (consultores):} mensalidade entre US\$ 30 e US\$ 50 para acesso completo.
    \item \textbf{Enterprise:} licença anual de US\$ 2.000 a US\$ 5.000 para uso local, customização e apoio especializado.
\end{enumerate}

Este modelo é consistente com a lógica de mercado regional: baixo custo de aquisição, foco em produtividade e uma arquitetura local que reduz a dependência de infraestrutura física e de cloud.

\section{Estimativa Económico-Financeira Preliminar}
\begin{table}[H]
\centering
\renewcommand{\arraystretch}{1.3}
\begin{tabularx}{\textwidth}{X|ccc}
\toprule
\textbf{Indicador} & \textbf{Ano 1} & \textbf{Ano 2} & \textbf{Ano 3} \\
\midrule
Subscrições Pro & 50 & 150 & 350 \\
Licenças Enterprise & 0 & 2 & 5 \\
Receita anual (USD) & 24.000 & 76.000 & 196.000 \\
Custos fixos (USD) & 5.000 & 12.000 & 25.000 \\
Resultado operacional (USD) & 19.000 & 64.000 & 171.000 \\
Ponto de equilíbrio & 8º mês & consolidado & consolidado \\
\bottomrule
\end{tabularx}
\caption{Projeção financeira preliminar para 3 anos.}
\end{table}

A premissa base assume entrada no mercado em 2027, com custos operacionais reduzidos e foco em crescimento orgânico e consultivo. A natureza \textit{asset-light} do projeto reforça a viabilidade financeira e neta simulação.

\section{Matriz de Riscos e Mitigação}
\begin{table}[H]
\centering
\renewcommand{\arraystretch}{1.2}
\begin{tabularx}{\textwidth}{>{\raggedright\arraybackslash}X|c|c|>{\raggedright\arraybackslash}X}
\toprule
\textbf{Risco} & \textbf{Prob.} & \textbf{Impacto} & \textbf{Mitigação} \\
\midrule
Resistência à mudança para Excel & Média & Alto & Interface amigável e importação direta de dados. \\
Custo de aquisição do cliente & Alto & Médio & Modelo Freemium + Pro incluso. \\
Concorrência de software enterprise & Média & Alto & Foco em acessibilidade e segurança local. \\
Adoção lenta no mercado empresarial & Média & Alto & Primeiro foco em consultores e estudantes. \\
\bottomrule
\end{tabularx}
\caption{Riscos estratégicos e operacionais.}
\end{table}

\section{Impacto Socioeconómico e Ambiental}
\begin{itemize}
    \item \textbf{Impacto socioeconómico:} criação de propriedade intelectual angolana, retenção de talento e fortalecimento da inovação local no setor energético.
    \item \textbf{Impacto ambiental:} melhoria da eficiência operacional, redução de desperdícios e aumento da precisão na gestão de reservatórios, com potencial para reduzir extração excessiva e otimizar decisões de desenvolvimento.
\end{itemize}

\section{Guião de Pitch de 5 Minutos (GEOENERGY)}
\begin{enumerate}[leftmargin=*]
    \item \textbf{Gancho (30s):} “Quantos de vocês já passaram horas a tentar calcular produção com Excel ou ficaram impossibilitados de utilizar software de reservatórios porque a licença custa mais do que o orçamento do projeto?”
    \item \textbf{Problema (45s):} “A engenharia de reservatórios continua dependente de ferramentas fragmentadas e caras. Pequenos consultores e estudantes não têm acesso profissional a soluções robustas, rápidas e certeiras.”
    \item \textbf{Solução (60s):} “O PetroLúmen é uma plataforma desktop que integra PVT, análise de declínio, simulação e IA em uma única interface. Funciona localmente, importa dados rapidamente e entrega decisões mais rápidas.”
    \item \textbf{Prova e modelo (60s):} “O repositório já inclui backend em Python/FastAPI, frontend em React e Tauri, e uma base de testes automatizados. O modelo de negócio parte por estudantes, segue para consultores e depois escala para empresas.”
    \item \textbf{Impacto e CTA (45s):} “O PetroLúmen é um projeto feito por angolanos, para o mercado global. Pedimos ao GEOENERGY mentoria, rede e apoio para transformar esta ideia numa solução exportável e relevante para a indústria energética.”
\end{enumerate}

\section{Conformidade com o Regulamento GEOENERGY}
A proposta cumpre a lógica do regulamento porque:
\begin{itemize}
    \item enquadra-se na categoria \textbf{Digitalização e Tecnologias Aplicadas às Geociências e Energia};
    \item tem equipa com até 5 elementos e um proponente principal identificado;
    \item descreve problema, solução, valor para o cliente, mercado, viabilidade técnica e potencial de impacto;
    \item segue a estrutura recomendada para o Modelo de Negócio (Anexo I)
    \item apresenta uma proposta em formato PDF, conforme a exigência de submissão.
\end{itemize}

\end{document}
